\documentclass[conference]{IEEEtran}
\IEEEoverridecommandlockouts

% Paquetes esenciales
\usepackage{cite}
\usepackage{amsmath,amssymb,amsfonts}
\usepackage{algorithmic}
\usepackage{graphicx}
\usepackage{textcomp}
\usepackage{xcolor}

\def\BibTeX{{\rm B\kern-.05em{\sc i\kern-.025em b}\kern-.08em
    T\kern-.1667em\lower.7ex\hbox{E}\kern-.125emX}}

\begin{document}

\title{Laboratory Report 1: Implementation of Logic Circuits}

% Autores separados correctamente usando \and
\author{
    \IEEEauthorblockN{Cristopher J. Balam-Góngora}
    \and
    \IEEEauthorblockN{Daniel A. Batun-Leon}
    \and
    \IEEEauthorblockN{Christian A. Cetina-Sanchez}
    \and
    \IEEEauthorblockN{Erick E. Canul-Segovia}
    \and
    \IEEEauthorblockN{Javier Rosado-Penes}
} 

\maketitle

\begin{abstract}
% Escribe tu resumen aquí (máximo 150-250 palabras).
\end{abstract}

\begin{IEEEkeywords}
% Escribe entre 3 y 5 palabras clave separadas por comas.
Logic circuits, Boolean algebra, gates
\end{IEEEkeywords}

\section{Introduction}
% Contexto de la práctica, objetivos e introducción teórica.

\section{Methodology}
% Materiales utilizados, diseño del circuito, diagramas.

\section{Results and Discussion}
% Tablas de verdad empíricas, mediciones, problemas encontrados y cómo se resolvieron.

\section{Conclusion}
% Conclusiones principales derivadas de los resultados obtenidos.

\section*{Acknowledgment}
% (Opcional) Agradecimientos a profesores o instituciones.

% Bibliografía
\begin{thebibliography}{00}

\bibitem{ref1} 
% Ejemplo de referencia
% A. B. Author, ``Title of paper,'' Journal Name, vol. x, pp. xxx--xxx, Year.

\end{thebibliography}

\end{document}